\documentclass[11pt,aps,preprint,nofootinbib,superscriptaddress,tightenlines]{revtex4-2}
\usepackage{amsmath,amssymb,bm,mathtools,graphicx}
\usepackage{xcolor}
\definecolor{darkblue}{RGB}{0,0,160}
\usepackage[colorlinks=true,linkcolor=darkblue,citecolor=darkblue,urlcolor=darkblue]{hyperref}

\def\st{\begin{equation}}
\def\stp{\end{equation}}
\def\Eq#1{eq.~(\ref{#1})}
\def\Sect#1{Sect.~\ref{#1}}
\newcommand{\avg}[1]{\langle #1\rangle}
\let\oldparagraph\paragraph
\renewcommand{\paragraph}{\par\bigskip\oldparagraph}

\begin{document}

\title{Starting in equilibrium: energy and momentum fluctuations at a global temperature}
\author{D.~Teaney}
\affiliation{Department of Physics and Astronomy, Stony Brook University,
Stony Brook, New York 11794, USA}
\date{\today}

\begin{abstract}
The present thermal start (\texttt{thermal\_init!}) draws the particle momenta at fixed
particle energies, so the energy fluctuations are missing, and the stochastic step,
which holds every particle energy fixed, never creates them.  Here every particle
energy and momentum is drawn from the equilibrium at one global $\beta^\mu$: the
maximum of the total entropy at fixed total energy and momentum, and the fluctuations
about it.  We define the global $\beta^\mu$ from the average totals, draw configurations
whose totals fluctuate with the thermodynamic susceptibilities, optionally shift them
to the exact targets so that a run needs no thermalization, and test the draws.
\end{abstract}

\maketitle

%=============================================================================
\section{Maximum entropy at fixed totals}
\label{sec:max}

At fixed positions each particle has a volume $V_a$, an energy $e_a$ and a momentum
$\bm m_a$, and its entropy is the function
\st
S_a(e_a,\bm m_a)=V_a\,\sigma\big(e_a/V_a,\ \bm m_a/V_a\big),
\stp
with $\sigma(\varepsilon,\bm M)$ the entropy density at energy density $\varepsilon=T^{00}$ and
momentum density $\bm M=T^{0i}$ (\texttt{stochastic.tex}).  The totals
\st
E=\sum_ae_a,\qquad \bm M_{\rm tot}=\sum_a\bm m_a
\stp
are conserved.  The most probable state maximizes $\sum_aS_a$ at fixed $E$ and
$\bm M_{\rm tot}$.  With multipliers $\beta^0$ and $\bm\beta$, each particle maximizes the same
function,
\st
\Phi_a(e_a,\bm m_a)=S_a(e_a,\bm m_a)-\beta^0e_a+\bm\beta\cdot\bm m_a ,
\label{Phi}
\stp
so that at the maximum
\st
\frac{\partial S_a}{\partial e_a}=\beta^0=\frac1T,\qquad
\frac{\partial S_a}{\partial\bm m_a}=-\bm\beta=-\frac{\bm v}{T} ,
\label{stat}
\stp
with $T=T_{\rm rel}/\gamma$ the bare temperature and $\bm v$ the fluid velocity, common
to all particles.  In covariant form, with $\beta_\mu=u_\mu/T_{\rm rel}$ and the metric
$(-,+,+)$, $\Phi_a=S_a+\beta_\mu T^{0\mu}_a$, since $\beta_0=-\gamma/T_{\rm rel}=-\beta^0$ and
$\beta_i=u^i/T_{\rm rel}=\beta^i$.  The particle quantity in \Eq{Phi} is the entropy
function $S_a$ of the particle's energy and momentum, not the entropy $s_a$ itself;
$s_a=S_a(e_a,\bm m_a)$ once $e_a$ and $\bm m_a$ are drawn.

%=============================================================================
\section{Fluctuations and the global $\beta^\mu$}
\label{sec:fluct}

\paragraph{One particle.}  In equilibrium the probability of a configuration is
$e^{\sum_aS_a}$ at the fixed totals.  For a system of many particles each particle is
then distributed independently as
\st
p_a(e,\bm m)\propto e^{\Phi_a(e,\bm m)}=e^{S_a(e,\bm m)-\beta^0e+\bm\beta\cdot\bm m} .
\label{pa}
\stp
To quadratic order about the maximum, $x=(\delta e,\delta\bm m)$ is Gaussian with covariance
\st
C_a=\big(-\partial^2S_a\big)^{-1},
\label{Ca}
\stp
the particle's susceptibility.  At rest, with $c_v=d\varepsilon/dT$ and $w=e+p$,
\st
\avg{\delta e^2}=V_aT^2c_v,\qquad \avg{\delta m^i\delta m^j}=V_a\,wT\,\delta^{ij} .
\stp
For the conformal equation of state $V_aT^2c_v=4V_a\varepsilon T$.

\paragraph{The average is not the maximum.}  The distribution \Eq{pa} is skewed: the mean
of $e$ lies above its most probable value by $0.16$ standard deviations at $s_a=720$
and by $0.58$ at $s_a=54$, and at $s_a=54$ the momentum variance is $16\%$ above
\Eq{Ca}.  We therefore define the global $\beta^\mu$ by the \emph{average} totals,
\st
\sum_a\avg{e_a}_\beta=E,\qquad \sum_a\avg{\bm m_a}_\beta=\bm M_{\rm tot} .
\label{avg}
\stp
With $\theta=(-\beta^0,\bm\beta)$ conjugate to $x=(e,\bm m)$, \Eq{pa} is an exponential family,
so $\partial\avg{x}/\partial\theta=\mathrm{Cov}(x)$ exactly, and \Eq{avg} is solved by Newton's
method,
\st
\theta\to\theta+\Big[\sum_a\mathrm{Cov}_a(x)\Big]^{-1}\Big[(E,\bm M_{\rm tot})-\sum_a\avg{x_a}\Big],
\stp
starting from the most probable state at the mean densities.  The resulting $T$ is
$0.5\%$ below the temperature of the most probable state at $s_a=541$ and $5\%$ below at
$s_a=54$.

\paragraph{Fluctuating or exact totals.}  The particles are drawn independently from
\Eq{pa}, the canonical ensemble; the totals then fluctuate about $(E,\bm M_{\rm tot})$
with the global susceptibility
\st
\mathrm{Cov}\big(E,\bm M_{\rm tot}\big)=\sum_a\mathrm{Cov}_a(x)\ \simeq\ V_{\rm tot}\,\chi ,
\label{global}
\stp
$V_{\rm tot}=\sum_aV_a$ and $\chi$ the susceptibility of $T^{0\mu}$ per volume.  To start an
SPH run at exactly the target $E$ and $\bm M_{\rm tot}$, the difference between the drawn
totals and the targets is then removed uniformly,
\st
x_a\to x_a-\frac{V_a}{V_{\rm tot}}\Big[\sum_bx_b-(E,\bm M_{\rm tot})\Big],
\stp
a uniform shift of the energy and momentum densities (an equal share per particle on an
equal-volume lattice).  This changes each particle's distribution only at order $1/N$,
so the run can start without a thermalization step.

%=============================================================================
\section{Implementation}
\label{sec:impl}

In \texttt{src/measure.jl}:
\begin{itemize}
\item \texttt{particle\_moments}: the exact $\avg{x}$ and $\mathrm{Cov}(x)$ of \Eq{pa} by
  quadrature on a $41^3$ grid around the maximum ($-6$ to $+10$ standard deviations in
  $e$, $\pm9$ in $\bm m$); particles of equal volume share them.
\item \texttt{beta\_from\_averages}: the Newton solution of \Eq{avg}.
\item \texttt{draw\_particle}: one exact draw from \Eq{pa}, an independence Metropolis
  chain of 30 steps whose proposal is the Gaussian of \Eq{Ca} with twice its covariance
  (wide enough to cover the skewed tail); the acceptance is $0.59$.
\item \texttt{equilibrium\_init!(ps, params, cl, eos; exact\_totals = false)}: from the
  totals of the current state (or given $E$, $\bm M_{\rm tot}$), solves for $\beta^\mu$,
  draws $(e_a,\bm m_a)$ at the current volumes $V_a=s_a/S_a$, sets $s_a=S_a(e_a,\bm m_a)$
  and re-smooths.  With \texttt{exact\_totals = true} the totals are then shifted to
  the targets as above.
\end{itemize}

%=============================================================================
\section{Tests}
\label{sec:tests}

\texttt{test/testeqstart.jl}; conformal equation of state, $h/a=2$, box $0.6$
($N=672$), $V_a=5.4\times10^{-4}$.

\paragraph{One particle.}  The draws of \texttt{draw\_particle} ($4\times10^4$) agree
with direct quadrature of \Eq{pa} in mean, variances and the $e$--$m_x$ covariance, at
rest and at $v=0.5$, at $s_a\simeq700$ and at $s_a\simeq50$, within the statistical
errors; the maximum satisfies \Eq{stat} to $10^{-6}$.

\paragraph{The lattice.}  With \texttt{exact\_totals = true} the totals are exact to $10^{-12}$,
the averages at the solved $\beta^\mu$ reproduce them, and the spread of $(e_a,\bm m_a)$
over the particles equals the single-particle covariance (at rest and at $v=0.32$).

\paragraph{Global susceptibility.}  $3000$ canonical configurations of the lattice:
\begin{center}
\begin{tabular}{lccccc}
\hline
 & $\avg{E}/E-1$ & $\frac{\mathrm{Var}\,E}{\sum_a C_a^{ee}}$ & $\frac{\mathrm{Var}\,M_x}{\sum_a C_a^{xx}}$ & corr$(E,M_x)$ & predicted \\ \hline
$s_a=541$, rest & $(-1.3\pm0.7)\times10^{-4}$ & $0.988$ & $0.968$ & $-0.003$ & $0$ \\
$s_a=541$, $v=0.32$ & $(-1.4\pm0.7)\times10^{-4}$ & $0.996$ & $0.975$ & $0.724$ & $0.727$ \\
$s_a=54$, rest & $(-2.0\pm2.0)\times10^{-4}$ & $0.967$ & $0.991$ & $-0.008$ & $0$ \\
\hline
\end{tabular}
\end{center}
(statistical error of the variance ratios $\pm0.026$).  At rest and $s_a=541$,
$\mathrm{Var}\,E=0.993\,V_{\rm tot}T^2c_v$ and $\mathrm{Var}\,M_x=0.983\,V_{\rm tot}wT$; at
$s_a=54$, $\mathrm{Var}\,M_x=1.17\,V_{\rm tot}wT$, the non-Gaussian excess, which \Eq{global}
with the exact $\mathrm{Cov}_a$ includes.  The two means at $s_a=541$ lie $2\sigma$ low with
the same sign, a possible bias of order $10^{-4}$ to be checked with more configurations.

%=============================================================================
\section{The SPH representation after the draw}
\label{sec:sph}

The draw holds the volumes $V_a$ fixed.  After re-smoothing, the code's volumes are
$V_a'=s_a/S_a'$, which follow the drawn entropy: a particle with more entropy has more
volume at nearly the same smoothed density.  Its energy recomputed from the SPH state,
$e_a'=s_a\,\mathcal E(S_a',\dots)$, then changes by
\st
\delta e_a'=\frac{\varepsilon}{S}\,\delta s_a=\tfrac34\,T\,\delta s_a
\qquad\text{instead of}\qquad \delta e_a=T\,\delta s_a
\stp
at fixed momentum (conformal, $\varepsilon=\tfrac34TS$), so $\mathrm{Var}\,e'/\mathrm{Var}\,e\simeq9/16$.
Measured: $0.62$, with $\mathrm{corr}(e',e)=0.997$ and $\mathrm{rms}(V'/V-1)=0.07$.

In a coupled run (leapfrog and stochastic step, $S=10^6$) the variance of the
recomputed particle energies, in units of the drawn one, stays at $0.57$--$0.61$ over
$t=3$ from this start, and at $0$--$0.004$ from \texttt{thermal\_init!}: on this time
scale the dynamics neither create nor remove these fluctuations, so they must be put in
at the start.  Which energy variance is the equilibrium one for the particle
representation, $\mathrm{Var}\,e$ or $\tfrac{9}{16}\mathrm{Var}\,e$, is open.

\end{document}
